\scanpage{99}

\deleted{If $N$ is \unsure{equivalent to the}} \deleted{Under these conditions,}

Then the function $u\to N(u)$ on $L^N_{\deleted{\unsure{0}}}(E,F)$ is a \uline{norm}. Indeed, we have $N(u+v)=\lim_{n\to\infty}N(|(\rho_i(u+v))|_n)\deleted{\unsure{,\rho_2(u+v)}}$, now the second side is bounded above by $N(|(\rho_i(u))|_n)\deleted{\unsure{,\rho_2(u)}}+N(|(\rho_i(v))|_n)\deleted{\unsure{,\rho_2(v)}}$ (number , formula (44)), hence by $N(u)+N(v)$. We also note trivially $N(\lambda u)=|\lambda|N(u)$, and $N(u)=0\Rightarrow u=0$.

\margin{$N(u)=N(u^*)$ (12)\par $\left.\begin{aligned}N(uv)&\leq\lVert u\rVert N(v)\\N(vu)&\leq\lVert u\rVert N(v)\end{aligned}\right\}$ (13)}

The space $L^N_{\deleted{\unsure{0}}}(E,F)$ is therefore a normed space; \deleted{\unsure{let us prove that it is \uline{complete}.}}\inserted{let} $L_N(E,F)$ be the normed subspace consisting of the $u$ such that $\unsure{(\rho_i(u))|(\rho_i u)}\in\ell_N$. I say that $L_N(E,F)$ is the closure of $\overline E\otimes F$ in $L^N(E,F)$. Indeed, let us show that

$\sum\rho_i\bar e_i\otimes f_i$ is the limit in the sense of $N$ of $\sum_1^n\rho_i\bar e_i\otimes f_i$, hence that $N(\sum_{n+1}^\infty\rho_i\bar e_i\otimes f_i)=N(\rho_{n+1}\unsure{(u)},\ldots)\to 0$ for $n\to+\infty$, in other words, that in $\ell_N$ the operations $x\to t_n(x)=(\overbrace{0{-}0}^{n},x_{n+1},x_{n+2}\unsure{-})$ tend to $0$ pointwise. Now they form an equicontinuous sequence, which tends to zero on \deleted{\unsure{a}} $R^{(N)}$, which proves the thing. \deleted{\unsure{Let us note} \unsure{therefore|$N_\infty$}, \unsure{when we have} $L^N(E,F)$ \unsure{and if|and that} $N$ \illegible{3}}

\margin{Definition of $L^N(E,F)$ if $N\simeq N_\infty$ ($L^N(E,F)=L(E,F)$, with \textcircled{14} $\boxed{N(u)=\underset{\substack{v\in F'\otimes E\\N^\circ(v)\leq 1}}{\operatorname{Sup}}\operatorname{Tr}(vu)}$ $L_N(E,F)$ as before). Case where $u$ is compact.}

\deleted{The dual of $L_N(E,F)$ is identified with a subspace of the dual of $L_{N_1}(E,F)=L^1(E,F)$, hence of $L(F,E)$, let us suppose that $N$ is not equivalent to $N_1$, hence $N'$ not equivalent to $N_\infty$. I say that}

\[
\deleted{(L_N(E,F))'=L^{N'}(F,E).}
\]

\deleted{In the first place, \deleted{\unsure{if}} if $u\in L_N(E,F)$, $v\in L^{N'}(F,G)$ then $uv\in L^1(E,G)$, with}

\[
\deleted{\lVert uv\rVert_1\leq N(u)N'(v)}
\]

\deleted{let} Let $N,N'\ N''$ be Schatten norms such that

\[
\boxed{N(xy)\leq N'(x)N''(y)}\tag*{(\unsure{15|13})}
\]

for \unsure{all|all} $x,y\in R^{(N)}$. Then the same formula will be valid for $x\deleted{\unsure{,y|y}}\in\ell^{N'}$, $y\in\ell^{N''}$. Let us then prove that $u\in L^{N'}(E,F)$, $v\in L^{N''}(F,G)$ implies $uv\in L^N(E,G)$, with

\[
\boxed{N(uv)\leq N'(u)N''(v)}\tag*{(16)}
\]

\margin{100}
